\section{Validation, deliverables, and the next decisive steps}
\label{sec:validation}

\subsection{Exact reference programs}

The Python programs use standard-library arbitrary integers and
\code{fractions.Fraction}; no floating-point decisions enter their
certificate checks or solution lists. The main checker expands the
unit and action residuals as sparse rational polynomials and compares
coefficients. It computes positive common denominators and absolute
column sums, enumerates the mixed-radix coordinate grid, and checks the
original equations at every retained point. Accepted example certificates
are exported as JSON files so the checked data can be inspected or
processed by another implementation.

The featured tests include a four-point elliptic section, rational
coordinates with distinct denominators, a quotient with nilpotents,
an empty rational zero set, a zero-radius grid, a unit expressed as a
linear combination rather than a literal generator, noncommuting
matrices on a redundant spanning family, and a localization that removes
an irrelevant positive-dimensional component. Independent linear-system
and factored-polynomial families compare against solutions known directly
from the input construction. Malformed and altered certificates are
rejected before enumeration. Grid decoding is separately compared with
a Cartesian-product implementation.

The separated-annihilator program also constructs multivariable monomial
rectangles and all action witnesses from the supplied ideal identities.
Its tests cover mixed degrees, degree one, tensor products, repeated roots
and the elliptic section. The norm program independently constructs monic reductions, polynomial
multiplication matrices and their determinants. Its determinant values
are compared with scalar Gaussian elimination at rational sample inputs.
It tests repeated and singular fibers, zero first coordinates, nonzero
constant norms, extraneous norm roots, the final $G=0$ filter, and
nontrivial rational denominators. Its norm-to-module constructor retains
the division quotients and adjugate witnesses. The resulting certificates
are checked by the separate coefficient checker, and its complete lists
are compared with the coordinate-grid enumeration.

These are implementation checks on finite cases. They do not establish
the universal theorems by testing. The universal mathematical arguments
are the proofs in Sections~\ref{sec:finite-observable-modules}
and~\ref{sec:norm-enumeration}; the exact tests detect mistakes in
the concrete programs and examples.

\subsection{Lean and mathematical review}

The standalone Lean build checks the eight modules listed in
Section~\ref{sec:formalization}. The archive records the compiler
version, successful module results, printed axiom dependencies and
the source-provenance check. The example module evaluates the elliptic
certificate independently: the checker returns true, the column bounds
are $[2,1]$, the grid has fifteen candidates, and the result has the
four displayed rational points. An evaluation diagnostic is kept distinct
from a theorem and from an object \ipra\ derivation.

Independent reviews checked the matrix orientation, use of the unit
certificate, denominator divisibility, maximum-coordinate estimate,
constructor residuals, elliptic example, localization identities,
the norm bridge, and the distinction between global and tuple-local
covers. The review records are included. They do not assert that all
earlier repository packets are correct. The PDF is compiled from the
included source and visually inspected after rendering; build and
layout results are recorded separately from mathematical validation.

\subsection{A useful integration order}

The first integration target is the proof-transport core: it is already
compiled and has a small dependency footprint. The next target is the
integer/rational finite arithmetic library described in the formalization
plan. Completing that library supports both enumeration methods and many
existing theorem packets.

For the mathematics, a particularly useful follow-up is a catalogue of
explicit finite sections and their certificates. Each entry should provide
the equations, the norm or action data, an independent identity check,
the returned list, and the precise role of that section in a larger
argument. This is more actionable than adding an abstract field saying
that a desired cover exists. The harder global theorem should then state
exactly which of those certificates it produces and how its coefficients
are bounded.

The current package completes a real interface improvement: finite
algebraic data previously used to count points now supply their
coordinates and complete lists, and a checked logical transformation
preserves local arithmetic proofs through finite computations. The
remaining global arithmetic and analytic obligations are exposed as
specific theorems, with no claim that changing the representation has
already proved them.
